% !PRINTED 438
\weberpage{438}
%<<CONT
\[
h = 2m,\ 4m,\ 6m,\ \ldots (n-1)m
\]
assumes, one obtains in (21) neither two equal nor two opposite values, and all the mutually distinct values of $x^2$ result.

Thus for an \gesperrt{odd} $n$ we have
\[
B_n(x) = \prod\limits_{1,\,\frac{n-1}{2}}^{\nu} \left(x^2 - 4\cos^2\frac{2\nu m\pi}{n}\right). \tag{24}
\]

If in this formula we set
\[
x = 2\cos\varphi,\quad x^2 = 4(1-\sin^2\varphi),\quad \cos^2\frac{2\nu m\pi}{n} = 1 - \sin^2\frac{2\nu m\pi}{n},
\]
we can represent it thus:
\[
\frac{\sin n\varphi}{\sin\varphi} = (-1)^{\frac{n-1}{2}} 2^{n-1} \prod\limits_{1,\,\frac{n-1}{2}}^{\nu} \left(\sin^2\varphi - \sin^2\frac{2\nu m\pi}{n}\right), \tag{25}
\]
from which again, for $\varphi=0$, formula (20) follows.

Since $\cos n\varphi$ vanishes for odd $n$ when
\[
\cos\varphi = \pm \sin\frac{2\nu m\pi}{n}
\]
holds, we have thereby found the decomposition into factors of the function $A_n(x)$ for this case also. For \gesperrt{odd} $n$ we have
\[
A_n(x) = x \prod\limits_{1,\,\frac{n-1}{2}}^{\nu} \left(x^2 - 4\sin^2\frac{2\nu m\pi}{n}\right). \tag{26}
\]

The quantities
\[
\alpha = \sin^2\frac{2\nu m\pi}{n} \tag{27}
\]
are thus the roots of an algebraic equation of degree $(n-1):2$
\[
\Phi_n(x) = 0,
\]
and we obtain the function $\Phi_n(x)$ if, in place of the variable $x^2$ in the function $A_n(x):x$, we put a new variable $4x$ and divide by $2^{n-1}$. It therefore becomes [formula (11)]
\[
\Phi_n(x) = \overset{\alpha}{\Pi}(x-\alpha)
= x^{\frac{n-1}{2}} - \frac{n}{4} x^{\frac{n-3}{2}} + \frac{n(n-3)}{1.2}\,\frac{1}{16}\, x^{\frac{n-5}{2}} - \cdots, \tag{28}
\]
and for (25) we can write
\[
\frac{\sin n\varphi}{\sin\varphi}
= (-1)^{\frac{n-1}{2}} 2^{n-1} \overset{\alpha}{\Pi}(\sin^2\varphi - \alpha). \tag{29}
\]
